%% kapa: exam=1 round=26 year=2015 period=1 form=A subject=economics subject_ko=경제원론 seq=35 number=75
%% source: Q-Net 감정평가사 기출문제 (cst00309) artlSeq=5138001 file=extracted/1008/5138001/제26회_감정평가사 1차_기출/제26회_감정평가사 1차_기출/1교시_감평_A형.pdf
%% conversion: kordoc 4.24.51 -> tools/gen1.py (auto); review=proofread-sonnet; flags=none
\begin{question}{75}
오쿤의 법칙(Okun's Law)에 따라 실업률이 1\% 포인트 증가하면 실질GDP는 약 2\% 포인트 감소한다고 가정하자. 만약, 중앙은행이 화폐공급 증가율을 낮추어 인플레이션율은 10\%에서 8\%로 하락하였으나 실업률은 4\%에서 8\%로 증가하였을 경우, 희생비율(sacrifice ratio)은?
(단, 희생비율$=\dfrac{\text{실질 GDP 감소율}}{\text{인플레이션 하락률}}$이다.)
\begin{choices}
  \item 약 2
  \item 약 4
  \item 약 6
  \item 약 8
  \item 약 10
\end{choices}
\end{question}
\answer{2}
